\section{Invariant resonance data versus representations}
\label{sec:unified-dictionary}

\subsection{The invariant core}
\label{subsec:invariant-core}
% BEGIN CURATED SUBJECT INDEX
\index{unbounded operator}
\index{Hilbert space}
\index{spectrum!point}
\index{resolvent}
\index{Jost function}
\index{Wronskian}
\index{connection coefficient}
\index{boundary condition!incoming and outgoing}
\index{resonance!residue}
% END CURATED SUBJECT INDEX

Across the constructions in the preceding Stages, the invariant core consists
of four items:
an analytically continued singularity, its spectral sheet, its residue or root
subspace, and the continuum contour relative to which it is exposed.  The
wave function, norm, and operator domain depend on the representation.  A
claim of equivalence should therefore compare resolvent matrix elements or
connection coefficients, not the pointwise form of two regularized wave
functions.

For a simple resonance, the comparison is most usefully read from the
connection problem outward:
\begin{align}
  \Cincoming(\resonantE,\hbar)=0
  &\Longleftrightarrow
  F(k_{\mathrm R})=0
  \notag\\
  &\Longleftrightarrow
  S(k)\text{ has a pole at }k_{\mathrm R}
  \notag\\
  &\Longleftrightarrow
  W(k_{\mathrm R})=0
  \notag\\
  &\Longleftrightarrow
  H_\theta\Psi_\theta
  =\resonantE\Psi_\theta .
  \label{eq:full-pole-chain}
\end{align}
Here $F$ is a Jost function, $W$ is an incoming/outgoing Wronskian, and
$\Cincoming$ is an exact-WKB connection coefficient.  The complex-scaled
eigenvalue relation holds for an angle that exposes the pole.  The rigged
Hilbert space supplies the pole functional associated with the common residue.

\subsection{A compact dictionary}
\label{subsec:compact-dictionary}
% BEGIN CURATED SUBJECT INDEX
\index{Hilbert space}
\index{locally convex space}
\index{topology}
\index{spectrum!point}
\index{resolvent}
\index{connection coefficient}
\index{boundary condition!incoming and outgoing}
\index{resonance!residue}
\index{lateral Borel sum}
\index{exact WKB}
% END CURATED SUBJECT INDEX

\begin{longtable}{>{\raggedright\arraybackslash}p{0.20\textwidth}
                  >{\raggedright\arraybackslash}p{0.34\textwidth}
                  >{\raggedright\arraybackslash}p{0.36\textwidth}}
\toprule
Language & Resonance object & Reliability test \\
\midrule
\endfirsthead
\toprule
Language & Resonance object & Reliability test \\
\midrule
\endhead
Exact WKB & Zero of a Borel-resummed incoming connection coefficient & Stokes
graph, lateral sum, sheets, and monodromy are documented \\
\addlinespace
Scattering theory & Pole of $S$, $T$, or a continued resolvent & Sheet,
outgoing boundary value, residue, and threshold structure are specified \\
\addlinespace
Complex scaling & Isolated eigenvalue of $H_\theta$ & Stability under
$\theta$ and basis changes; separation from the rotated continuum \\
\addlinespace
Feshbach projection & Nonlinear root of
$\det[z-\vsub{H}{eff}(z)]$ & Retarded branch of the self-energy and its
energy dependence are retained \\
\addlinespace
Rigged Hilbert space & Pole functional in $\PhiSpace^\times$ & Test space and
dual topology admit the required analytic continuation \\
\addlinespace
Zel'dovich method & Gaussian-regularized bilinear residue & The regulator is
removed by analytic continuation, not by an ordinary real-axis limit \\
\addlinespace
Black-hole perturbation & Pole of the retarded radial Green function & Horizon
and outer boundary conditions match the time convention; cuts and residues
are included \\
\bottomrule
\end{longtable}

The table is deliberately phrased in terms of tests rather than slogans.  For
example, ``complex eigenvalue'' is absent from the first column because it is
not by itself a definition.  Likewise, ``normalizable state'' means different
things before and after deformation.

\subsection{What non-self-adjointness does and does not mean}
\label{subsec:meaning-nonselfadjoint}
% BEGIN CURATED SUBJECT INDEX
\index{operator!adjoint and self-adjoint}
\index{spectrum!point}
\index{Riemann sheet}
\index{complex scaling}
\index{Feshbach projection}
% END CURATED SUBJECT INDEX

A non-self-adjoint operator may arise in at least four ways.  It may be a
fundamental effective generator after environmental degrees of freedom are
discarded.  It may be an energy-dependent Feshbach operator.  It may be an
analytic representation, such as $H_\theta$, of a self-adjoint scattering
problem.  Or it may be a finite numerical approximation whose
non-self-adjointness includes discretization artifacts.  Physical conclusions
depend on which situation applies.

Complex scaling belongs primarily to the third category.  It does not claim
that microscopic quantum mechanics violates unitarity.  Instead, it moves a
pole from a nonphysical sheet into the point spectrum of a deformed operator.
The rotated continuum completes the representation and restores the analytic
information needed to reconstruct real-energy unitary scattering.

An open-system Lindblad generator belongs to the first category after a
Markovian approximation.  Its eigenvalues describe decay of density-matrix
modes, not automatically poles of a one-particle $S$ matrix.  Quantum-jump
unravelings can relate the two in particular models, but complex scaling should
not be identified with a no-jump Hamiltonian without an explicit derivation
\cite{Piilo:2008dmi,Breuer2016NonMarkovian,Gneiting:2020pew}.

\subsection{Observables beyond pole positions}
\label{subsec:beyond-poles}
% BEGIN CURATED SUBJECT INDEX
\index{spectrum!point}
\index{resolvent}
\index{Wigner time delay}
\index{connection coefficient}
\index{resonance!residue}
\index{exact WKB}
\index{complex scaling}
\index{exceptional point}
\index{holonomy}
\index{quasinormal mode!excitation factor}
% END CURATED SUBJECT INDEX

A mature resonance calculation should aim at more than a table of complex
energies.  The next layer consists of residues, channel projectors, phase
shifts, time delays, and continuum level densities.  For black holes, this
means excitation factors, source-dependent Green functions, greybody factors,
and branch-cut discontinuities.  For exceptional points, it means the full
Jordan resolvent and a normalized holonomy, not only the eigenvalue square
root.

Exact WKB and complex scaling are complementary at this layer.  A derivative
of the exact connection coefficient gives a pole residue, whereas a
left--right complex-scaled eigenvector gives the same residue in an $L^2$
representation.  Establishing this equality numerically in nontrivial models
would be a stringent test of the unified framework.
